\section{Mapa de aplicaciones: cómo los flujos de trabajo digitales se convierten en el principal campo de batalla del agente}

\subsection{Forma de la tarea: de la operación en una sola aplicación a la orquestación entre aplicaciones}
El docente enfatiza que las tareas reales de los usuarios normalmente no son una página con un solo clic, sino un proceso que atraviesa varias herramientas. Por ejemplo: buscar información en el navegador, organizar datos en una hoja de cálculo, enviar los resultados por correo electrónico, archivar el registro en un documento. Un agente multimodal necesita manejar un estado persistente y el contexto de la tarea, no solo responder a instrucciones de corto plazo.

\begin{importantbox}{Cuatro características de las tareas del mundo real}
\begin{enumerate}
    \item \textbf{Cadena larga}: a menudo supera los 20--50 pasos;
    \item \textbf{Entre interfaces}: navegador, escritorio, terminal y herramientas de Office aparecen mezclados;
    \item \textbf{No determinismo}: la carga de páginas, las ventanas emergentes y las fluctuaciones de red provocan cambios de estado;
    \item \textbf{Objetivo abierto}: muchas tareas admiten varias rutas viables, no una única respuesta estándar.
\end{enumerate}
\end{importantbox}

\subsection{Estratificación de las capacidades del agente: percepción, planeación, ejecución, verificación}
El curso descompone al agente en una pila de capacidades, en lugar de tratarlo como un modelo de caja negra. La capa de percepción se encarga de entender la interfaz y los objetos semánticos; la capa de planeación decide el siguiente subobjetivo; la capa de ejecución produce acciones ejecutables; la capa de verificación determina si la tarea se completó y corrige la trayectoria. Esta estratificación permite ubicar mejor los cuellos de botella del sistema.

\begin{knowledgebox}{Por qué el agente multimodal necesita ``visibilidad del proceso''}
Cuando una tarea falla, el equipo debe saber si el problema está en la percepción, la planeación o la ejecución. Si todos los fallos se atribuyen a que ``el modelo no es lo bastante fuerte'', no es posible formular una estrategia de optimización eficaz. La observabilidad es el núcleo de la eficiencia de iteración del agente.
\end{knowledgebox}

\subsection{Prioridad a los escenarios digitales, los escenarios físicos después}
Desde el punto de vista del ritmo de implementación, el docente coloca claramente los entornos digitales (web, escritorio, móvil) como la dirección de mayor prioridad. La razón es que el ciclo de retroalimentación es más rápido, los guiones de evaluación son más fáciles de construir y el costo es más controlable. Los agentes en el mundo físico son importantes, pero su infraestructura de datos y evaluación es más pesada y su velocidad de iteración es más lenta.

\begin{warningbox}{No confundir prematuramente la ruta de la robótica con la del agente de GUI}
Ambos son agentes, pero sus restricciones de ingeniería son completamente distintas. El cuello de botella del agente de GUI está en la comprensión de la interfaz y la orquestación de tareas; el agente físico añade además el error de percepción, las restricciones de seguridad y la estabilidad del actuador. Discutirlos juntos oculta la naturaleza del problema.
\end{warningbox}

\subsection{Resumen de la sección}
El curso define el problema de las aplicaciones como una tarea de largo plazo que es ``entre aplicaciones, de múltiples pasos y recuperable''. De ahí se desprende el requisito de ingeniería: primero cerrar a fondo el ciclo de las tareas digitales, y después expandirse a escenarios más complejos.

